\chapter{Классификация банковских транзакций}

\section{Постановка задачи и выбор метода}

Банковская выписка в сыром виде представляет собой неразмеченный корпус коротких текстов: каждая строка содержит дату, сумму, направление платежа и произвольно заполненное поле назначения, вроде «Оплата по счёту 1254 от ООО Прогресс» или «Пл.\ по дог.\ б/о». Извлечь из этого управленческую информацию --- структуру расходов, долю возвратов, динамику выручки --- без семантической разметки невозможно. Задача сводится к автоматической категоризации коротких текстов по набору экономически значимых классов~\cite{aggarwal2012survey}.

В классической постановке задача категоризации текстов решается методами машинного обучения: наивным байесовским классификатором, методом опорных векторов на TF-IDF признаках~\cite{sebastiani2002machine} или предобученными трансформерами типа BERT~\cite{devlin2018bert}. Перечисленные подходы предполагают наличие размеченной обучающей выборки --- корпуса текстов с правильными метками, на котором обучается классификатор. В условиях малого бизнеса такой выборки нет изначально: ни одна транзакция не несёт готовой метки категории, а ручная разметка 4\,086 операций потребовала бы значительных ресурсов до начала полезного анализа. Кроме того, финансовый аудит требует \textit{объяснимости} каждого решения: финансовый менеджер должен понимать, почему операция отнесена в ту или иную категорию, и при необходимости исправить результат~\cite{whittington2012principles}.

Этим требованиям отвечает \textbf{гибридный подход}: rule-based классификатор как основной метод, большая языковая модель (LLM) как инструмент обработки неоднозначных назначений и ручная валидация выборки для оценки итогового качества. Rule-based системы исторически доминировали в задачах информационного извлечения до распространения статистических методов и сохраняют высокую эффективность там, где паттерны предметной области хорошо структурированы~\cite{Cohen1995}. В банковской выписке таких паттернов большинство: налоговые платежи, кредитные выплаты и коммунальные расходы имеют почти стандартизированные назначения. LLM привлекается именно для оставшегося «хвоста» --- операций с нестандартными или неполными описаниями, где детерминированные правила не срабатывают.

\section{Таксономия категорий}

Проектирование таксономии категорий --- самостоятельная исследовательская задача, предшествующая любой автоматической классификации. Без чёткого и непротиворечивого справочника классификатор не имеет устойчивого целевого пространства, а результаты разных запусков могут быть несопоставимы~\cite{pustejovsky2012natural}.

Для данного проекта разработан справочник из \textbf{17 укрупнённых категорий}, объединяющих все типы операций малого производственного предприятия. При проектировании выдержаны три методологических принципа. Во-первых, \textit{взаимная исключительность}: каждая операция относится ровно к одной категории, смысловые пересечения исключены. Это требование принципиально для агрегации в отчёт о движении денежных средств --- любое двойное отнесение искажало бы итоговые суммы~\cite{ias7}. Во-вторых, \textit{тип-зависимость}: для каждой категории задан допустимый тип операции (\texttt{credit} для поступлений, \texttt{debit} для списаний). Нарушение этого ограничения --- например, отнесение исходящего платежа в категорию «Выручка от клиентов» --- выявляется автоматически ещё на этапе классификации. В-третьих, \textit{управленческая интерпретируемость}: каждая категория напрямую соответствует строке управленческого ДДС согласно принципам МСФО~\cite{ias7}, что позволяет без промежуточных преобразований агрегировать транзакции в финансовую отчётность.

\begin{table}[H]
\centering
\caption{Таксономия категорий банковских операций}
\label{tab:categories}
\begin{tabular}{llc}
\toprule
\textbf{Категория} & \textbf{Секция ДДС} & \textbf{Тип} \\
\midrule
Выручка от клиентов & Операционный приток & credit \\
Прочие поступления & Операционный приток & credit \\
Закупки для производства & Операционный отток & debit \\
Упаковка и расходники & Операционный отток & debit \\
Персонал & Операционный отток & debit \\
Аренда и коммунальные & Операционный отток & debit \\
Связь, сервисы, подписки & Операционный отток & debit \\
Логистика и доставка & Операционный отток & debit \\
Маркетинг и реклама & Операционный отток & debit \\
Банковские комиссии & Операционный отток & debit \\
Профессиональные услуги & Операционный отток & debit \\
Налоги, взносы, штрафы & Операционный отток & debit \\
Возвраты & Корректировка & credit \\
Кредиты и займы & Финансовая деятельность & credit / debit \\
Внутренние переводы & Нейтрально & credit / debit \\
Вывод собственников & Финансовая деятельность & debit \\
Неизвестно & --- & credit / debit \\
\bottomrule
\end{tabular}
\end{table}

Отдельного внимания заслуживает категория «Возвраты»: в производственном бизнесе, работающем через дискаунтеры, возвраты нереализованной продукции представляют собой существенную корректировку выручки и не могут объединяться с «Выручкой от клиентов» без искажения реальной картины продаж. Выделение возвратов в самостоятельную категорию позволяет рассчитать брутто-выручку, нетто-выручку и процент возврата по каждому товару независимо.

\section{Rule-based классификатор}

Классический rule-based подход к категоризации текстов предполагает набор правил вида «если выполняется условие $C_i$, то присвоить метку $L_i$», применяемых в порядке убывания специфичности~\cite{Cohen1995}. Конкурентное преимущество этого подхода в финансовом контексте --- полная интерпретируемость: для каждой транзакции можно указать конкретное правило, сработавшее при классификации, и при необходимости скорректировать его. Это критически важно при аудите, где каждое решение должно быть верифицируемым~\cite{whittington2012principles}.

В проекте реализована пятиуровневая иерархия правил с убыванием детерминированности. \textbf{Первый уровень} --- детерминированные идентификаторы по ИНН: налоговая инспекция и государственные фонды (ПФР, СФР, ФСС) имеют уникальные ИНН, полностью определяющие категорию вне зависимости от содержания поля назначения. Этот уровень обрабатывает все налоговые и страховые платежи с уверенностью 1,0. \textbf{Второй уровень} --- семантические паттерны в назначении платежа, реализованные как регулярные выражения над нормализованным текстом. Примеры: шаблон \texttt{/зарплат|аванс|подотчёт|отпускн/i} идентифицирует выплаты персоналу; \texttt{/мука|сахар|меланж|яичн|масло.раст/i} --- сырьевые закупки; \texttt{/упаковк|пакет|плёнк|стрейч/i} --- упаковочные материалы. \textbf{Третий уровень} --- контрагентный словарь: справочник юридических лиц, для которых категория известна заранее (например, ООО «РАДУГА Н» $\to$ закупки для производства, ИП Баскаков С.А.\ $\to$ аренда). \textbf{Четвёртый уровень} --- широкие правила для типичных паттернов без точного совпадения (поступления от юридических лиц $\to$ выручка). \textbf{Пятый уровень} --- присвоение категории «Неизвестно» операциям, не подпавшим ни под одно правило.

Каждому правилу присвоен скалярный показатель уверенности: \texttt{high} ($\geq 0{,}90$) для детерминированных идентификаторов и точных совпадений, \texttt{medium} ($0{,}70$--$0{,}89$) для семантических паттернов с умеренной специфичностью, \texttt{low} ($< 0{,}70$) для широких правил. Показатель уверенности используется в двух целях: для выбора между конкурирующими классификаторами (rule-based vs LLM) и для флагирования операций, требующих ручной проверки. Операция помечается флагом \texttt{needs\_review} при уверенности ниже 0,70, при категории «Неизвестно» или при попадании суммы в топ-5\% по абсолютному значению. Последнее условие охватывает крупные нетиповые платежи, которые непропорционально влияют на финансовый результат даже при правильной автоматической классификации.

\section{Большие языковые модели для разметки транзакций}

Большие языковые модели (LLM) продемонстрировали высокую точность в задачах классификации коротких текстов с нестандартными формулировками благодаря обширным языковым знаниям, полученным при предобучении~\cite{brown2020language}. В банковском контексте LLM способна интерпретировать деловой жаргон, неполные аббревиатуры и нестандартные описания, которые оставляют rule-based классификатор без ответа.

Ключевым ограничением LLM при работе с финансовыми данными является конфиденциальность: передача банковских выписок во внешние API-сервисы создаёт юридические и репутационные риски. В проекте используется \textbf{локально развёрнутый LLM-сервер} (LM Studio) с открытой языковой моделью, что полностью исключает передачу данных за пределы рабочей машины. Компромисс --- несколько более низкое качество по сравнению с коммерческими frontier-моделями --- приемлем, поскольку LLM используется лишь для «хвоста» неоднозначных операций, а не для основной массы.

Второй принципиальный риск LLM --- галлюцинации: модель может возвращать категорию, не входящую в допустимый справочник, или нарушать ограничения по типу операции~\cite{ji2023survey}. Для предотвращения этого промпт явно перечисляет все 17 допустимых категорий, указывает запрещённые комбинации (расходные категории недопустимы при \texttt{credit}, доходные --- при \texttt{debit}), и требует ответа исключительно в формате JSON с полями \texttt{category}, \texttt{subcategory} и \texttt{confidence}. Ответы, не прошедшие валидацию по схеме, отбраковываются и соответствующие операции получают флаг \texttt{needs\_review}.

Необходимо также отметить существенное ограничение предложенного подхода. Использование LLM для классификации банковских транзакций --- даже в локальном развёртывании --- несёт значительные инфраструктурные издержки: установку и обслуживание LM Studio, выбор и периодическое обновление модели, мониторинг качества ответов и валидацию JSON-схемы. Для предприятия малого бизнеса с небольшим числом контрагентов этот инструмент является избыточным: хорошо спроектированный rule-based классификатор с контрагентным словарём закрывает 90--95\% операций с высокой уверенностью~\cite{Cohen1995}. LLM оправдывает себя там, где объём «неоднозначного хвоста» велик, паттерны назначений разнообразны и нет ресурсов на поддержку обширного словаря. В данном проекте LLM в итоге применена лишь к единственной операции из 4\,086 --- результат, который заставляет поставить под сомнение целесообразность её включения в производственный контур без существенного масштабирования задачи~\cite{sculley2015hidden}.

Финальная категория каждой операции (\texttt{category\_final}) определяется по схеме приоритетов: при уверенности rule-based классификатора $\geq 0{,}90$ используется его результат; при уверенности ниже порога --- результат LLM, если он прошёл валидацию; при расхождении между двумя классификаторами --- операция направляется на ручную проверку. Этот каскадный алгоритм проиллюстрирован на рисунке~\ref{fig:classification_flow}.

\begin{figure}[H]
\centering
\begin{tikzpicture}[
  node distance=1.1cm and 2.8cm,
  every node/.style={font=\small},
  process/.style={rectangle, rounded corners=5pt, draw=black!60, fill=blue!10,
    text width=5.8cm, minimum height=0.95cm, align=center},
  decision/.style={diamond, aspect=2.5, draw=black!60, fill=orange!12,
    text width=3cm, align=center, inner sep=0pt},
  result/.style={rectangle, rounded corners=5pt, draw=green!60!black, fill=green!10,
    text width=4.2cm, minimum height=0.85cm, align=center},
  alert/.style={rectangle, rounded corners=5pt, draw=red!60, fill=red!8,
    text width=4.2cm, minimum height=0.85cm, align=center},
  arrow/.style={-Stealth, thick, black!65}
]
  \node[process] (input) {Банковская операция\\(direction, purpose, INN)};
  \node[process, below=of input] (rule) {Rule-based классификатор\\(5 уровней иерархии)};
  \node[decision, below=of rule] (conf) {confidence\\$\geq 0{,}90$?};
  \node[result, right=2.8cm of conf] (rule_final) {\textbf{category\_final = rule}\\(94,2\% операций)};
  \node[process, below=1.5cm of conf] (llm) {LLM-классификатор\\(LM Studio, JSON output)};
  \node[decision, below=of llm] (agree) {rule = llm\\или llm valid?};
  \node[result, right=2.8cm of agree] (llm_final) {\textbf{category\_final = llm}\\(прочие операции)};
  \node[alert, below=1.5cm of agree] (review) {\textbf{needs\_review = True}\\238 операций (5,8\%)};

  \draw[arrow] (input) -- (rule);
  \draw[arrow] (rule) -- (conf);
  \draw[arrow] (conf) -- node[above]{\small Да} (rule_final);
  \draw[arrow] (conf) -- node[right]{\small Нет} (llm);
  \draw[arrow] (llm) -- (agree);
  \draw[arrow] (agree) -- node[above]{\small Да} (llm_final);
  \draw[arrow] (agree) -- node[right]{\small Нет} (review);
\end{tikzpicture}
\caption{Каскадная схема классификации банковских операций}
\label{fig:classification_flow}
\end{figure}

\section{Результаты классификации и оценка качества}

По результатам применения гибридного классификатора к 4\,086 банковским операциям достигнута практически полная разметка датасета (таблица~\ref{tab:classification_results}). Единственная операция, оставшаяся в категории «Неизвестно», составляет 5\,793 руб. --- 0,01\% от совокупного оборота. 238 операций (5,8\% по количеству) направлены на ручную проверку как потенциально спорные.

\begin{table}[H]
\centering
\caption{Итоги классификации банковских операций}
\label{tab:classification_results}
\begin{tabular}{lr}
\toprule
\textbf{Показатель} & \textbf{Значение} \\
\midrule
Всего операций & 4\,086 \\
Классифицировано rule-based (conf $\geq 0{,}90$) & 3\,847 (94,2\%) \\
Классифицировано через LLM & 1 (0,0\%) \\
Направлено на ручную проверку & 238 (5,8\%) \\
Осталось unknown & 1 (0,01\% оборота) \\
Сумма операций на проверке & 31\,380\,000 руб. \\
\bottomrule
\end{tabular}
\end{table}

Принципиально важное наблюдение: на ручную проверку направлено лишь 5,8\% операций по количеству, но их совокупная сумма --- 31,4 млн руб. Это объясняется тем, что флаг \texttt{needs\_review} срабатывает в том числе для крупных нетиповых платежей (топ-5\% по сумме). Таким образом, 94,2\% рутинных операций обработаны автоматически, а внимание финансового менеджера сконцентрировано на тех операциях, которые несут наибольший управленческий вес.

Итоговое распределение расходов по категориям представлено в таблице~\ref{tab:category_breakdown} и на рисунке~\ref{fig:expenses_bar}. Структура расходов подтверждает корректность разметки: доминирование производственных закупок (69,4\% расходов) соответствует ожидаемому профилю пищевого производства~\cite{sme_profitability_ru}, а суммы по категориям «Аренда» и «Персонал» соответствуют данным управленческого учёта.

\begin{table}[H]
\centering
\caption{Расходы по категориям (итоговая классификация, май 2025 --- апрель 2026)}
\label{tab:category_breakdown}
\begin{tabular}{lrrr}
\toprule
\textbf{Категория} & \textbf{Операций} & \textbf{Сумма, руб.} & \textbf{Контраг.} \\
\midrule
Закупки для производства & 321 & 29\,804\,203 & 14 \\
Аренда и коммунальные & 56 & 4\,642\,629 & 6 \\
Упаковка и расходники & 50 & 3\,252\,206 & 5 \\
Персонал & 106 & 2\,971\,955 & 6 \\
Налоги, взносы, штрафы & 66 & 1\,156\,005 & 2 \\
Профессиональные услуги & 36 & 783\,501 & 3 \\
Кредиты и займы & 3 & 234\,480 & 1 \\
Связь, сервисы и подписки & 41 & 144\,464 & 4 \\
Вывод собственников & 1 & 35\,000 & 1 \\
\bottomrule
\end{tabular}
\end{table}

\begin{figure}[H]
\centering
\begin{tikzpicture}
\begin{axis}[
  xbar,
  width=13cm, height=6.5cm,
  bar width=0.42cm,
  xlabel={Сумма, млн руб.},
  ytick=data,
  yticklabels={
    Прочие,
    Связь и сервисы,
    Банковские комиссии,
    Профессиональные услуги,
    Налоги и взносы,
    Персонал,
    Упаковка,
    Аренда и ЖКХ,
    Закупки},
  yticklabel style={font=\small},
  xmin=0, xmax=33,
  xmajorgrids=true,
  grid style={dashed,gray!40},
  axis y line=left,
  axis x line=bottom,
  nodes near coords,
  nodes near coords align={horizontal},
  every node near coord/.append style={font=\scriptsize},
]
\addplot[fill=blue!55, draw=blue!70!black] coordinates {
  (0.035,1) (0.144,2) (0.099,3) (0.784,4)
  (1.156,5) (2.972,6) (3.222,7) (4.643,8) (29.161,9)
};
\end{axis}
\end{tikzpicture}
\caption{Структура операционных расходов по категориям (млн руб.), май 2025 --- апрель 2026}
\label{fig:expenses_bar}
\end{figure}

Для оценки качества автоматической разметки подготовлена стратифицированная валидационная выборка из 500 операций, включающая все 238 операций с флагом \texttt{needs\_review}~\cite{pustejovsky2012natural}. По результатам ручной разметки рассчитываются метрики согласованности: accuracy и macro-F1 для каждого из трёх классификаторов (rule-based, LLM, финальный). Такая методология, основанная на выборочной проверке вместо полной ручной разметки, является стандартной практикой при оценке систем автоматической аннотации в отсутствие заранее размеченного золотого стандарта~\cite{pustejovsky2012natural}.

Итог этапа: 4\,086 банковских операций переведены из неструктурированных текстовых строк в размеченный датасет с категориями, соответствующими строкам управленческого ДДС. Это создаёт основу для построения финансовой отчётности и обучения прогнозных моделей в последующих главах.
